\chapter{Clearness}

This was not Tibor McMasters’s first encounter with runners. In Charlottesville they came and went freely, tolerated by all. Wherever runners lived, a peculiar calm prevailed—a local peace generated not by authority but by the runners’ own mild, ingratiating nature.

They gathered around him now, peering up with bright, friendly curiosity. They were small, scarcely four feet high, rotund and furred, with twitching noses, beady eyes, and powerful hind legs like kangaroos. They spoke together, as if rehearsed.

“Clearness be with you,” they said. “Why don’t you have any arms or legs? You are very strange.”

“The war,” Tibor replied shortly.

They examined his cart with professional concern.

“Your wheel bearings are failing,” one said.

Alarm shot through him. He lifted a wheel and spun it. The dry grinding confirmed their diagnosis.

“There is an autofac nearby,” the largest runner said. “It still works a little. It could repair your bearings.”

Relief flooded Tibor. He agreed at once. The runners scampered ahead, chanting gleefully about guarantees lasting a thousand years or a million miles.

Before following, Tibor turned to the lizards. “Can I trust them?” he asked.

Potter assured him that runners were harmless. “Unlike bugs,” he added, kicking one aside. The insect scuttled away.

They reached the autofac at dawn. The sky blazed with infantile colors; mutant birds chirped among weeds. Earl, leader of the runners, hesitated, trying to remember the exact approach. Finally he found the entrance: a bare circular clearing, dominated by a sealed metal disc embedded in the earth.

A Russian autofac, Tibor thought. Seeded from orbit near the end of the war.

He greeted it politely.

Earl recoiled in alarm. “Don’t speak casually to it. It can kill us.”

Irritated, Tibor persisted. He asked for help plainly, refusing ritual obeisance. Earl, groaning, performed the proper invocation instead, begging the autofac’s mercy.

The machine responded erratically.

A bullhorn emerged and accused Tibor of pregnancy, offering grotesque folk remedies. When corrected, the autofac grew hostile, vented gas, stuttered incoherently, and finally collapsed into silence amid smoke and mechanical shrieks.

“I think you killed it,” a runner said mournfully.

Earl looked stricken. “It probably couldn’t have helped anyway.”

Tibor urged the cow onward. Earl made one last attempt, scribbling a written request and lowering it into the entrance, but Tibor did not wait. His cart clattered painfully as he left, the bearings nearly seized.

As he traveled, a strange melancholy settled over him—not despair, but calm. He sang fragments of hymns, mangling the words, comforted nonetheless. Then the grinding stopped.

The wheel had frozen entirely.

He halted the cow. This, he knew, was as far as he could go.

The morning world stirred around him. Birds sang. Life—damaged, grotesque, persistent—went on. One bird sang words, fragments of hymns, rearranged and reassembled.

A meta-mutant, Tibor thought. A forward oddity.

Before he could reflect further, the bushes to his right exploded outward.

A colossal worm emerged.

It uncoiled, shrieking, dragging itself on oily slime. It proclaimed its immortality, its origin in waste and war, its guardianship of something precious. It spat, raged, demanded that Tibor leave.

“I can’t,” Tibor said.

The worm attacked. Slime coated the cart and cow. Tibor activated the cart’s electric field; sparks crackled. The worm recoiled, then lunged again.

He fired.

Twice.

The worm collapsed, wounded, accusing him.

“I’m sorry,” Tibor said, reloading. “Only one of us could live.”

He fired again.

The worm died.

Tibor examined the slime, considering whether it could lubricate his bearings. Then he recalled the worm’s words—my possessions—and guided the cart past the corpse.

Behind the bushes lay a cave filled with junk: rusted machines, broken weapons, old appliances, posters, magazines. Worthless hoarded debris. The worm had guarded nothing.

Disappointment washed over him.

The hymn-singing bird fluttered closer.

“You can understand me now,” it said.

“Yes,” Tibor replied. “I can.”

“You touched the worm,” the bird explained. “Now you hear all of us.”

The bird knew his name. Knew his purpose.

“You are searching for the God of Wrath,” it said. “To paint him.”

Tibor’s chest tightened.

“You know where he is?” he asked.

“Yes,” the bird replied calmly. “He is not far from here. I can lead you to him—if you wish.”
